%% ====================================================================
%%  Chapter 3 — Methodology
%% ====================================================================
\chapter{Methodology}
\label{ch:methodology}

This chapter describes the system, data, models, and evaluation methods employed in this thesis. First,
we review the DLoc system that introduced the CSI-to-heatmap problem. Second, we describe the WILD
dataset on which our experiments are based. Third, we describe the four models evaluated in this thesis.
Finally, we describe the experiments performed to evaluate each of the models.

\section{Overview of the DLoc System Architecture}
\label{sec:method:dloc}

The DLoc system~\cite{ayyalasomayajula2020deep} frames the indoor localization problem as an
image-to-image translation problem. The image provided to the system is built from the CSI measurements
collected from the WiFi access points, and the system outputs a heatmap that represents the location of
the device emitting the WiFi signal. More specifically, each of the four access points provides a
two-dimensional representation of the CSI measurements collected from that access point, and these
representations are stacked together into a four-channel image; the heatmap output by the system has the
same spatial dimensions as the input image but only a single channel, and contains a Gaussian bump
indicating the location of the device at a spatial resolution of $0.1$\,m per pixel.

DLoc's architecture consists of a ResNet encoder~\cite{he2016resnet} and a location decoder. The ResNet
encoder contains six residual blocks with 64 filters each, and compresses the four-channel image into a
high-dimensional representation. The location decoder then upsamples this representation to the size of
the original image to generate the heatmap. Additionally, DLoc employs a second decoder that outputs an
``offset'' heatmap to ensure that the outputs of the four access points are consistent with the location
of the device being localized. The DLoc model has an approximate parameter count of 16.5~million
parameters. The location of the device is determined by applying a sigmoid function to the output of the
location decoder and finding the peak of the resulting heatmap.

The training loss of DLoc is defined as a combination of the difference between the predicted and target
heatmaps and the offset heatmap, and the Adam optimizer ($\beta_1 = 0.5$) is used to minimize this objective. The
DLoc system is utilized in this thesis in two ways: as a baseline model against which the other models
are compared in terms of both accuracy and efficiency, and as the teacher in the training of the
knowledge-distilled student models.

\section{The WILD Dataset}
\label{sec:method:dataset}

All experiments are performed upon the WILD (WiFi-based Indoor Localization Dataset) dataset, which was
collected in the Jacobs Hall building at UC San Diego~\cite{ayyalasomayajula2020deep}. The dataset was
collected from four access points during five different data-collection sessions. Each collection
session took place under slightly different environmental conditions, as summarized in
Table~\ref{tab:method:sessions}.

More specifically, Sessions 1 and 2 (named \texttt{jacobs\_Jul28} and \texttt{jacobs\_Jul28\_2})
occurred within one hour of each other and used the same environmental conditions. During training, the
training data from Session 1 and the training data from Session 2 were both used to train the model;
during testing, the test data from Session 2 was used to evaluate the model. Sessions 3, 4, and 5
occurred three weeks later, during which the rooms where the measurements were collected were altered in
arbitrary ways (Session 3), given a different furniture layout (Session 4), or fitted with an aluminum
panel used to reflect the WiFi signal (Session 5).

\begin{table}[H]
\centering
\caption{The five WILD collection sessions used in this thesis.}
\label{tab:method:sessions}
\begin{tabular}{@{}llrl@{}}
\toprule
\textbf{Session} & \textbf{Tag} & \textbf{Samples} & \textbf{Condition} \\
\midrule
1 & \texttt{jacobs\_Jul28}        & ---     & training (reference conditions) \\
2 & \texttt{jacobs\_Jul28\_2}     & 4{,}260$^\dagger$ & in-distribution test (1\,h later, identical) \\
3 & \texttt{jacobs\_Aug16\_1}     & 11{,}201 & furniture rearrangement \\
4 & \texttt{jacobs\_Aug16\_3}     & 8{,}754  & different furniture configuration \\
5 & \texttt{jacobs\_Aug16\_4\_ref}& 5{,}486  & aluminum reflector (specular multipath) \\
\bottomrule
\end{tabular}

\vspace{0.4em}
{\footnotesize $^\dagger$ Held-out test split; Sessions~1+2 together provide 17{,}574 training
samples.}
\end{table}

\section{Proposed Student Architectures}
\label{sec:method:students}

The four models evaluated in this thesis have the same input and output specifications as the DLoc
model. Each model takes as its input the four-channel image generated from the CSI measurements of the
WiFi access points, and outputs a heatmap of the same spatial dimensions as the input image but with a
single channel, where a Gaussian bump at the peak indicates the location of the device. The four
students are summarized in Table~\ref{tab:method:architectures}.

\begin{table}[H]
\centering
\footnotesize % Minimized the text size for a smaller footprint
\setlength{\tabcolsep}{4pt} % Minimized the column spacing to tighten the table
\caption{Summary of the proposed student architectures and the DLoc teacher.}
\label{tab:method:architectures}
\begin{tabular}{@{}lrcl@{}}
\toprule
\textbf{Model} & \textbf{Parameters} & \textbf{Compression} & \textbf{Architecture highlights} \\
\midrule
DLoc baseline (teacher) & 16.5M & $1\times$ & ResNet encoder + location + offset decoders \\
MobileNetV2 & 2.3M & $7.2\times$ & depthwise-separable, inverted residuals \\
TinyCNN & 47K & $351\times$ & 3 conv + 2 FC layers \\
MobileNetV2-UNet & 2.41M & $6.8\times$ & MobileNetV2 encoder + U-Net skip connections \\
Mamba & 0.65M & $25\times$ & selective state-space sequence model \\
\bottomrule
\end{tabular}
\end{table}

\subsection{MobileNetV2 Student}
\label{sec:method:mobilenet}

The MobileNetV2 student~\cite{sandler2018mobilenetv2} incorporates the MobileNetV2 backbone in place of
the ResNet-based encoder of the DLoc model. The MobileNetV2 model's first convolutional layer is adapted
to accept the four channels of input from the WiFi access points. A decoder head then maps the features
extracted by the MobileNetV2 backbone to the heatmap output. This model contains approximately
2.3~million parameters.

\subsection{TinyCNN}
\label{sec:method:tinycnn}

TinyCNN is a minimal deep learning model that contains only three convolutional layers and two fully
connected layers. It contains approximately 47{,}000 parameters. This minimal model is included to
evaluate the limits of possible model sizes for this type of problem.

\subsection{MobileNetV2-UNet}
\label{sec:method:unet}

The MobileNetV2-UNet student~\cite{ronneberger2015unet} incorporates the same MobileNetV2 backbone as the
MobileNetV2 student, but in addition includes skip connections between the MobileNetV2 encoder and
decoder. The encoder is divided into four separate stages, each with a different number of channels; the
decoder concatenates the features from those encoder stages with the features from the decoder's
upsampled representations to create a heatmap. The first stage concatenates a 320-channel representation
with a 96-channel representation, the second combines a 128-channel representation with a 32-channel
representation, the third combines a 64-channel representation with a 24-channel representation, and the
final stage uses only the 32-channel representation. The MobileNetV2-UNet student contains 2.41~million
parameters, and is illustrated in Figure~\ref{fig:method:unet_arch}.

These skip connections are included because they allow the model to access the fine spatial features of
the WiFi signal that are otherwise lost when the representation of that signal is compressed to the
low-dimensional latent representation employed by the MobileNetV2 model.

\begin{figure}[H]
\centering
\resizebox{\linewidth}{!}{% ============================================================================
%  MobileNetV2-UNet architecture — \input-able body (no \documentclass).
%  Requires in the master preamble: \usepackage{tikz} and
%  \usetikzlibrary{positioning, arrows.meta, calc}.  Colours are defined here
%  so the figure is self-contained.  Layer specs from code/models/student_unet.py.
% ============================================================================
\definecolor{cEnc}{RGB}{142,202,230}
\definecolor{cBottle}{RGB}{33,102,172}
\definecolor{cDec}{RGB}{69,123,157}
\definecolor{cHead}{RGB}{188,108,37}
\definecolor{cSkip}{RGB}{42,157,143}
\definecolor{cIO}{RGB}{230,57,70}

\begin{tikzpicture}[
  font=\small,
  >={Stealth[length=2.4mm]},
  block/.style={rectangle, rounded corners=2pt, draw=black!70, thick,
                minimum height=11mm, minimum width=13mm, align=center,
                text=white, font=\footnotesize\bfseries},
  enc/.style={block, fill=cEnc, text=black},
  bottle/.style={block, fill=cBottle},
  dec/.style={block, fill=cDec},
  head/.style={block, fill=cHead},
  io/.style={rectangle, rounded corners=1pt, draw=black!70, thick,
             minimum height=13mm, minimum width=12mm, align=center,
             fill=cIO!12, font=\footnotesize},
  flow/.style={->, black!75, thick},
  skipline/.style={->, cSkip, very thick, dashed},
]
\node[io] (in) {Input\\$4{\times}161{\times}361$};
\node[enc,    right=7mm of in]  (inc) {inc\\$4{\to}32$\\/2};
\node[enc,    right=6mm of inc] (e1)  {enc1\\$24$};
\node[enc,    right=6mm of e1]  (e2)  {enc2\\$32$};
\node[enc,    right=6mm of e2]  (e3)  {enc3\\$96$};
\node[bottle, right=6mm of e3]  (e4)  {enc4\\$320$};
\node[dec,  right=11mm of e4] (d1) {dec1\\$416{\to}128$};
\node[dec,  right=6mm of d1]  (d2) {dec2\\$160{\to}64$};
\node[dec,  right=6mm of d2]  (d3) {dec3\\$88{\to}32$};
\node[head, right=6mm of d3]  (hd) {head\\$32{\to}1$\\Sigmoid};
\node[io,   right=7mm of hd]  (out) {Output\\$1{\times}161{\times}361$};
\draw[flow] (in)--(inc); \draw[flow] (inc)--(e1); \draw[flow] (e1)--(e2);
\draw[flow] (e2)--(e3);  \draw[flow] (e3)--(e4);  \draw[flow] (e4)--(d1);
\draw[flow] (d1)--(d2);  \draw[flow] (d2)--(d3);  \draw[flow] (d3)--(hd);
\draw[flow] (hd)--(out);
\draw[skipline] (e3.north) to[out=90,in=90]
      node[above,font=\tiny,text=cSkip]{skip 96} (d1.north);
\draw[skipline] (e2.north) to[out=90,in=90]
      node[above,font=\tiny,text=cSkip]{skip 32} (d2.north);
\draw[skipline] (e1.north) to[out=90,in=90]
      node[above,font=\tiny,text=cSkip]{skip 24} (d3.north);
\node[text=black!60, font=\scriptsize, below=3mm of e2.south] {MobileNetV2 encoder (pretrained, 4-ch input)};
\node[text=black!60, font=\scriptsize, below=3mm of d2.south] {U-Net decoder (upsample $+$ skip concat)};
\end{tikzpicture}
}
\caption{The MobileNetV2-UNet student. A MobileNetV2 encoder (light blue) downsamples the four-channel
CSI input; a decoder (dark blue) upsamples back to full resolution; teal dashed arcs are the U-Net skip
connections that concatenate encoder feature maps into the decoder at matching resolutions, bypassing the
bottleneck (deep blue). The head applies a $1{\times}1$ convolution and a sigmoid. Layer specifications
are exact, taken from the implementation.}
\label{fig:method:unet_arch}
\end{figure}

\subsection{State-Space (Mamba) Variant}
\label{sec:method:mamba}

The last model evaluated in this thesis is based on the Mamba selective state-space
model~\cite{gu2023mamba}. Whereas DLoc, the MobileNetV2 student, and the MobileNetV2-UNet student all
treat the CSI input as a two-dimensional image, the Mamba model treats the CSI measurements from each
access point as a sequence. The Mamba student contains only approximately 0.65~million parameters, making
it the most compressed of the models evaluated in this thesis. It is evaluated as a model with a
different inductive bias from the other students, to test whether this bias is competitive with those
other approaches to the CSI-to-heatmap problem.

\section{Knowledge-Distillation Framework}
\label{sec:method:distillation}

Each student was trained in two modes: a standalone mode that used only the ground-truth heatmaps for
supervision, and a distilled mode that additionally aimed to minimize the difference between its outputs
and the teacher model's output. The objective function used for the distilled mode is a convex
combination of the ground-truth heatmap term and the teacher-output term,
\begin{equation}
  \mathcal{L} \;=\; \alpha \cdot \mathrm{MSE}(\hat{y},\, y) \;+\; \beta \cdot \mathrm{MSE}(\hat{y},\, y_T),
  \label{eq:method:distill}
\end{equation}
where $\hat{y}$ is the student's output, $y$ is the ground-truth heatmap, and $y_T$ is the teacher's
precomputed output. The weights used in the combination are $\alpha = 0.7$ and $\beta = 0.3$. Because the
teacher is fixed, its output is precomputed prior to training, so the teacher's forward pass is never
invoked during the training of the students. Furthermore, because the teacher's targets are heatmaps
rather than one-hot maps, its outputs are inherently ``soft'': the probability mass is spread around the
true location rather than concentrated at a single point. These soft targets therefore encourage the
students to reproduce the location of the target device as well as the probability mass around it, which
acts as an implicit regularizer for the student models.

\section{Experimental Design}
\label{sec:method:design}

The experiments performed on the proposed models are described below.

\begin{enumerate}[leftmargin=2em]
  \item In-distribution accuracy is measured on the test set from Session 2, to determine the accuracy of
  the localization models and the cost of model compression.

  \item The models are evaluated under various forms of signal degradation at test time, to determine
  their robustness to real-world signal conditions. The corruptions applied include individual and random
  AP dropout, Gaussian noise of $\sigma \in \{0.1, 0.2\}$, global and single-AP attenuation, and spatial
  blur.

  \item Cross-session generalization is evaluated within the original DLoc protocol (each fold trains on the two July sessions plus two of the
  three August sessions, and tests on the held-out August session) as well as through a zero-shot evaluation in which the
  student models are trained only on the two July sessions and evaluated on the August sessions.

  \item The DLoc model is additionally trained using three different random seeds, to evaluate the
  statistical significance of the reported results.

  \item Lastly, the computational efficiency of each model is measured, an essential consideration for
  potential deployment of the methods.
\end{enumerate}

\section{Evaluation Metrics}
\label{sec:method:metrics}

The primary evaluation metric is the localization error of the models, which is the distance between the
predicted location (the peak of the predicted heatmap) and the location of the ground-truth heatmap. The
distance between the two points is given by
\begin{equation}
  e \;=\; 0.1 \times \big\lVert \arg\max(\hat{y}) - \arg\max(y) \big\rVert_2 ,
  \label{eq:method:error}
\end{equation}
scaled by the $0.1$\,m-per-pixel resolution of the CSI measurements. From the localization error for each
sample, the median error, 90th-percentile error (P90), 99th-percentile error (P99), and mean error are
calculated. Because the error is a scaled pixel distance, the achievable error values are quantized to
multiples of $0.1\sqrt{n}$; as a result, two models can have identical median errors yet different full
error distributions. In these cases the mean error and the error distribution itself are additionally
reported.